% GENERATED FILE. Edit canonical lessons, then run npm run generate:books.
% canonical-source-hash: fnv1a64:97b8b67e44abb51d
% canonical-lessons: SA-C395-sarthah, SA-C395-sarthavahah, SA-C395-anuyati, SA-C395-aprcchate, SA-C395-parivrajakah

\chapter{Caravans and Companions}
\label{ch:sa-caravans-and-companions}
\hlchaptermodality{395}

\begin{chapteropening}
\textbf{By the end of this chapter:} \emph{I can say who I travel with and how I take leave.}

Complete the last lesson of chapter 395: I can say who I travel with and how I take leave.
\end{chapteropening}

\section[\texorpdfstring{\sk{सार्थः} (sārthaḥ)}{सार्थः (sārthaḥ)}]{\sk{सार्थः} (sārthaḥ) — a caravan}

\label{lesson:SA-C395-sarthah}



\begin{quote}
Before the new one: say the Sanskrit for travels, then the Sanskrit for an ascetic's water-pot.
\end{quote}

\subsection*{You'll want to know: \sk{सार्थः}}
\textbf{\sk{सार्थः}} — \emph{sārthaḥ} — \textquotedblleft{}a caravan\textquotedblright{}.

\begin{quote}
\sk{सार्थः} \sk{ग्रामात्} \sk{ग्रामं} \sk{गच्छति।}
\end{quote}

\begin{quote}
\emph{sārthaḥ grāmāt grāmaṃ gacchati.}
\end{quote}

\begin{quote}
The caravan goes from village to village.
\end{quote}

A masculine noun.

\begin{cousinweb}
\textbf{\sk{सार्थः}} (\emph{sārthaḥ}) is built from \emph{sa}, \textquotedblleft{}with\textquotedblright{}, and \emph{artha}, \textquotedblleft{}an aim, business\textquotedblright{}: MW gives \emph{sārtha} first as having a purpose, then as a travelling company of traders or pilgrims, a caravan (MW).
\end{cousinweb}

\begin{grammarlens}[title={What you've built}]
One more word for getting around.
\end{grammarlens}

\subsection*{Your turn}
\begin{itemize}
\raggedright
  \item \emph{Say it:} \emph{sārthaḥ}
  \item \emph{Say it:} \emph{sārthaḥ}, once more
  \item \emph{Say it:} \emph{sārthaḥ grāmāt grāmaṃ gacchati.} — The caravan goes from village to village.
  \item \emph{Recall:} say the Sanskrit for travels, then the Sanskrit for an ascetic's water-pot, then say \emph{sārthaḥ} again
\end{itemize}

\begin{tcolorbox}[breakable,colback=teal!4,colframe=teal!35!black,title={Before you move on}]
What does \sk{सार्थः} mean? (\textquotedblleft{}A caravan\textquotedblright{}.) Say it once more.
\end{tcolorbox}
\section[\texorpdfstring{\sk{सार्थवाहः} (sārthavāhaḥ)}{सार्थवाहः (sārthavāhaḥ)}]{\sk{सार्थवाहः} (sārthavāhaḥ) — a caravan leader}

\label{lesson:SA-C395-sarthavahah}



\begin{quote}
Before the new one: say the Sanskrit for an ascetic's water-pot, then the Sanskrit for a caravan.
\end{quote}

\subsection*{You'll want to know: \sk{सार्थवाहः}}
\textbf{\sk{सार्थवाहः}} — \emph{sārthavāhaḥ} — \textquotedblleft{}a caravan leader\textquotedblright{}.

\begin{quote}
\sk{सार्थवाहः} \sk{मार्गं} \sk{सम्यक्} \sk{जानाति।}
\end{quote}

\begin{quote}
\emph{sārthavāhaḥ mārgaṃ samyak jānāti.}
\end{quote}

\begin{quote}
The caravan leader knows the road well.
\end{quote}

A masculine noun.

\begin{cousinweb}
\textbf{\sk{सार्थवाहः}} (\emph{sārthavāhaḥ}) is built from \emph{sārtha}, \textquotedblleft{}a caravan\textquotedblright{}, and \emph{vāha}, \textquotedblleft{}bearing, leading\textquotedblright{}, from the root √vah, \textquotedblleft{}to carry\textquotedblright{}; MW gives the compound as the leader of a caravan, a merchant (MW).
\end{cousinweb}

\begin{grammarlens}[title={What you've built}]
One more word for getting around.
\end{grammarlens}

\subsection*{Your turn}
\begin{itemize}
\raggedright
  \item \emph{Say it:} \emph{sārthavāhaḥ}
  \item \emph{Say it:} \emph{sārthavāhaḥ}, once more
  \item \emph{Say it:} \emph{sārthavāhaḥ mārgaṃ samyak jānāti.} — The caravan leader knows the road well.
  \item \emph{Recall:} say the Sanskrit for an ascetic's water-pot, then the Sanskrit for a caravan, then say \emph{sārthavāhaḥ} again
\end{itemize}

\begin{tcolorbox}[breakable,colback=teal!4,colframe=teal!35!black,title={Before you move on}]
What does \sk{सार्थवाहः} mean? (\textquotedblleft{}A caravan leader\textquotedblright{}.) Say it once more.
\end{tcolorbox}
\section[\texorpdfstring{\sk{अनुयाति} (anuyāti)}{अनुयाति (anuyāti)}]{\sk{अनुयाति} (anuyāti) — goes along with; follows}

\label{lesson:SA-C395-anuyati}



\begin{quote}
Before the new one: say the Sanskrit for a caravan, then the Sanskrit for a caravan leader.
\end{quote}

\subsection*{You'll want to know: \sk{अनुयाति}}
\textbf{\sk{अनुयाति}} — \emph{anuyāti} — \textquotedblleft{}goes along with; follows\textquotedblright{}.

\begin{quote}
\sk{मम} \sk{मित्रं} \sk{माम्} \sk{अनुयाति।}
\end{quote}

\begin{quote}
\emph{mama mitraṃ mām anuyāti.}
\end{quote}

\begin{quote}
My friend goes along with me.
\end{quote}

Past (he …): \sk{सः} \sk{माम्} \sk{अनुयातवान्}, \emph{saḥ mām anuyātavān}.

\begin{cousinweb}
\textbf{\sk{अनुयाति}} (\emph{anuyāti}) is built from \emph{anu}, \textquotedblleft{}after, along\textquotedblright{}, and the root √yā, \textquotedblleft{}to go\textquotedblright{}; MW gives \emph{anu-yā} as to go after or follow, with the participle \emph{anuyāta} (MW).
\end{cousinweb}

\begin{grammarlens}[title={What you've built}]
One more word for getting around.
\end{grammarlens}

\subsection*{Your turn}
\begin{itemize}
\raggedright
  \item \emph{Say it:} \emph{anuyāti}
  \item \emph{Say it:} \emph{anuyāti}, once more
  \item \emph{Say it:} \emph{mama mitraṃ mām anuyāti.} — My friend goes along with me.
  \item \emph{Recall:} say the Sanskrit for a caravan, then the Sanskrit for a caravan leader, then say \emph{anuyāti} again
\end{itemize}

\begin{tcolorbox}[breakable,colback=teal!4,colframe=teal!35!black,title={Before you move on}]
What does \sk{अनुयाति} mean? (\textquotedblleft{}Goes along with; follows\textquotedblright{}.) Say it once more.
\end{tcolorbox}
\section[\texorpdfstring{\sk{आपृच्छते} (āpṛcchate)}{आपृच्छते (āpṛcchate)}]{\sk{आपृच्छते} (āpṛcchate) — takes leave; says goodbye}

\label{lesson:SA-C395-aprcchate}



\begin{quote}
Before the new one: say the Sanskrit for a caravan leader, then the Sanskrit for goes along with.
\end{quote}

\subsection*{You'll want to know: \sk{आपृच्छते}}
\textbf{\sk{आपृच्छते}} — \emph{āpṛcchate} — \textquotedblleft{}takes leave; says goodbye\textquotedblright{}.

\begin{quote}
\sk{यात्री} \sk{मातरम्} \sk{आपृच्छते।}
\end{quote}

\begin{quote}
\emph{yātrī mātaram āpṛcchate.}
\end{quote}

\begin{quote}
The traveller takes leave of his mother.
\end{quote}

Past (he …): \sk{सः} \sk{मातरम्} \sk{आपृष्टवान्}, \emph{saḥ mātaram āpṛṣṭavān}.

\begin{cousinweb}
\textbf{\sk{आपृच्छते}} (\emph{āpṛcchate}) is built from \emph{ā}, \textquotedblleft{}towards\textquotedblright{}, and the root √prach, \textquotedblleft{}to ask\textquotedblright{}, which also gives \emph{pṛcchati}, \textquotedblleft{}asks\textquotedblright{}; in the middle voice MW gives \emph{ā-prach} as to take leave, to bid farewell (MW).
\end{cousinweb}

\begin{grammarlens}[title={What you've built}]
One more word for getting around.
\end{grammarlens}

\subsection*{Your turn}
\begin{itemize}
\raggedright
  \item \emph{Say it:} \emph{āpṛcchate}
  \item \emph{Say it:} \emph{āpṛcchate}, once more
  \item \emph{Say it:} \emph{yātrī mātaram āpṛcchate.} — The traveller takes leave of his mother.
  \item \emph{Recall:} say the Sanskrit for a caravan leader, then the Sanskrit for goes along with, then say \emph{āpṛcchate} again
\end{itemize}

\begin{tcolorbox}[breakable,colback=teal!4,colframe=teal!35!black,title={Before you move on}]
What does \sk{आपृच्छते} mean? (\textquotedblleft{}Takes leave; says goodbye\textquotedblright{}.) Say it once more.
\end{tcolorbox}
\section[\texorpdfstring{\sk{परिव्राजकः} (parivrājakaḥ)}{परिव्राजकः (parivrājakaḥ)}]{\sk{परिव्राजकः} (parivrājakaḥ) — a wandering ascetic}

\label{lesson:SA-C395-parivrajakah}



\begin{quote}
Before the new one: say the Sanskrit for goes along with, then the Sanskrit for takes leave.
\end{quote}

\subsection*{You'll want to know: \sk{परिव्राजकः}}
\textbf{\sk{परिव्राजकः}} — \emph{parivrājakaḥ} — \textquotedblleft{}a wandering ascetic\textquotedblright{}.

\begin{quote}
\sk{परिव्राजकः} \sk{ग्रामे} \sk{एकां} \sk{रात्रिं} \sk{वसति।}
\end{quote}

\begin{quote}
\emph{parivrājakaḥ grāme ekāṃ rātriṃ vasati.}
\end{quote}

\begin{quote}
The wandering ascetic stays one night in a village.
\end{quote}

A masculine noun.

\begin{cousinweb}
\textbf{\sk{परिव्राजकः}} (\emph{parivrājakaḥ}) is built from \emph{pari}, \textquotedblleft{}around\textquotedblright{}, and the root √vraj, \textquotedblleft{}to walk, to go\textquotedblright{}, with the agent ending \emph{-aka}; MW gives \emph{parivrājaka} as a wandering religious mendicant (MW).
\end{cousinweb}

\begin{grammarlens}[title={What you've built}]
One more word for getting around.
\end{grammarlens}

\subsection*{Your turn}
\begin{itemize}
\raggedright
  \item \emph{Say it:} \emph{parivrājakaḥ}
  \item \emph{Say it:} \emph{parivrājakaḥ}, once more
  \item \emph{Say it:} \emph{parivrājakaḥ grāme ekāṃ rātriṃ vasati.} — The wandering ascetic stays one night in a village.
  \item \emph{Recall:} say the Sanskrit for goes along with, then the Sanskrit for takes leave, then say \emph{parivrājakaḥ} again
\end{itemize}

\begin{tcolorbox}[breakable,colback=teal!4,colframe=teal!35!black,title={Before you move on}]
What does \sk{परिव्राजकः} mean? (\textquotedblleft{}A wandering ascetic\textquotedblright{}.) Say it once more.
\end{tcolorbox}
